\section{Speech and Audio Foundation Models: The Acoustic Toolkit}
\label{sec:foundationaudio}

Section~\ref{sec:humanmt} traced the translation stack whose \emph{source} is
already human language. A bark is not. To turn a continuous, noisy,
variable-length vocalization into something a translator could consume, one
needs the acoustic front-end and the audio-token machinery that the speech
community built for itself. This section catalogs that toolkit from a
computer-science standpoint---self-supervised encoders, recognition
architectures, the corpora and benchmarks that made scale possible, neural
audio codecs, and audio--language models---and asks, for each, what an animal
decoder inherits. The recurring lesson is that \emph{the engineering of the
front-end is largely solved; the missing piece is semantic, not acoustic}.

\subsection{Self-supervised speech representation learning}
The decisive idea is to learn a general acoustic representation from raw,
unlabeled waveform, then fine-tune on a small labeled task. Contrastive
Predictive Coding (CPC) showed that predicting future latent frames yields
transferable features~\cite{vandenoord2018cpc}; wav2vec extended this with a
contrastive objective over a discrete vocabulary~\cite{schneider2019wav2vec};
and wav2vec~2.0 combined a convolutional feature encoder, a quantized latent
space, and a masked Transformer to reach strong recognition with only minutes
of transcribed audio~\cite{baevski2020wav2vec2}. HuBERT replaced the
contrastive target with masked prediction of offline-clustered
units~\cite{hsu2021hubert}, WavLM added denoising and speaker-aware
pretraining~\cite{chen2022wavlm}, and data2vec unified the masked-prediction
objective across speech, vision, and text~\cite{baevski2022data2vec}, with a
faster successor in data2vec~2~\cite{baevski2023data2vec2}. Two extensions
matter directly for animals: wav2vec-U learns to recognize speech
\emph{without any transcriptions} by adversarially aligning discrete units to
a language model~\cite{baevski2021wav2vecu}, and multilingual scaling
(XLS-R, MMS, USM) pushes a single encoder across dozens to thousands of
languages~\cite{babu2022xlsr,pratap2023mms,zhang2023usm}. The animal analogue
is exact: a large unlabeled corpus of barks, pretrained with a masked or
contrastive objective, is the correct starting point, and the
transcription-free recipe is the closest human precedent for a target with no
parallel corpus.

\subsection{Recognition architectures and the alignment problem}
Before Transformers, end-to-end recognition was built on recurrent networks
with two alignment regimes. Connectionist Temporal Classification (CTC)
marginalizes over all alignments between input frames and output
labels~\cite{graves2006ctc}, and deep recurrent CTC systems scaled to
end-to-end speech recognition~\cite{graves2013speech,amodei2016deepspeech2};
attention-based listen-attend-spell models instead learn a soft
alignment~\cite{chan2016las}. The Conformer married convolution with
self-attention to capture both local and global acoustic
structure~\cite{gulati2020conformer}. For animal signals the alignment
problem is the same in form---a variable-length vocalization must be mapped
to a variable-length label sequence---but harder in substance, because there
is no canonical label sequence to align to (Section~\ref{sec:eval}).

\subsection{Corpora and benchmarks that made scale possible}
None of the above would exist without shared data and evaluation. LibriSpeech
provided a thousand hours of read audiobook
speech~\cite{panayotov2015librispeech}; Common Voice crowdsourced a massively
multilingual corpus~\cite{ardila2020commonvoice}; Multilingual LibriSpeech
extended read speech across languages~\cite{pratap2020mls}; and the older
TIMIT and Switchboard corpora anchored phonetic and conversational
research~\cite{garofolo1993timit,godfrey1992switchboard}. Benchmarks then
standardized what ``good'' means: SUPERB probes a frozen representation across
many speech tasks~\cite{yang2021superb}, while CoVoST and CoVoST~2 supply
speech-translation references~\cite{wang2020covost,wang2021covost2}. The
bioacoustic community is only now building the equivalent infrastructure
(BEANS, BirdSet; Section~\ref{sec:crossspecies}), and the absence of a
``LibriSpeech of barks'' is a concrete, fixable bottleneck.

\subsection{Neural audio codecs and discrete acoustic tokens}
A second lineage turns waveform into \emph{discrete tokens} that a language
model can predict. WaveNet modeled raw audio autoregressively~\cite{vandenoord2016wavenet};
VQ-VAE introduced a learned discrete latent codebook~\cite{vandenoord2017vqvae};
Jukebox showed that hierarchical VQ-VAEs could generate coherent
music~\cite{dhariwal2020jukebox}. SoundStream and EnCodec made neural audio
compression practical with residual vector
quantization~\cite{zeghidour2021soundstream,defossez2022encodec}. AudioLM then
demonstrated that a language model over these tokens can generate long-range
coherent audio without any text~\cite{borsos2023audiolm}, and the recipe
generalized to music (MusicLM, MusicGen) and general sound
(AudioGen)~\cite{agostinelli2023musiclm,copet2023musicgen,kreuk2023audiogen}.
VALL-E treated text-to-speech as neural codec language
modeling~\cite{wang2023valle}; UniAudio pushed toward a single foundation model
for many audio tasks~\cite{yang2024uniaudio}; and AudioPaLM, SALMONN, and
Qwen-Audio fused speech/audio encoders with large language
models~\cite{rubenstein2023audiopalm,tang2024salmonn,chu2024qwenaudio}. This is
precisely the discretization step a ``bark alphabet'' would require: if
vocalizations can be compressed into reusable tokens, sequence models can
search for structure---the strategy Project CETI applied to whale
codas~\cite{sharma2024whale}.

\subsection{The visual backbone analogy}
The same representation-learning revolution happened in vision, and the
analogy clarifies what transfers. AlexNet's ImageNet win established deep
convolutional features as general-purpose~\cite{krizhevsky2012alexnet,deng2009imagenet,russakovsky2015imagenet};
residual networks made very deep models trainable~\cite{he2016resnet}; and the
Vision Transformer later showed the Transformer backbone works on pixels as
well as tokens~\cite{dosovitskiy2021vit}. DogFACS-style facial analysis
(Section~\ref{sec:biology}) rides directly on this lineage. The lesson is that
a modality becomes tractable once a large pretrained backbone exists---and
that the backbone, not the task head, is where the value lives.

\subsection{Bioacoustic foundation models}
Finally, the toolkit has been pointed at animals. BirdNET demonstrated
large-scale, consumer-grade bird identification from
audio~\cite{kahl2021birdnet}; comprehensive reviews chart both the promise and
the methodological pitfalls of computational
bioacoustics~\cite{stowell2022review,stowell2014automatic}. Self-supervised
AVES~\cite{hagiwara2023aves}, the audio--language NatureLM-audio
model~\cite{robinson2024naturelm}, Google's Perch embeddings~\cite{ghani2023perch},
the BEANS benchmark~\cite{hagiwara2023beans}, and open general-purpose audio
encoders such as OpenBEATs~\cite{openbeats2025} together show that the
human speech recipe---pretrain on raw audio, transfer to many tasks---works
across species. What none of them supplies is a \emph{reference translation}:
they classify, embed, and retrieve, but they do not, and cannot yet, render a
bark as a sentence with measured fidelity. Table~\ref{tab:audio} condenses the
toolkit and its transfer value.

\begin{table}[htbp]
\centering
\caption{The speech/audio foundation-model toolkit and its transfer value to
animal decoding. ``Transfer value'' is our assessment. SSL = self-supervised;
ASR = automatic speech recognition; codec = discrete acoustic tokens.}
\label{tab:audio}
\small
\setlength{\tabcolsep}{4pt}
\begin{tabular}{@{}p{3.0cm}p{1.8cm}p{3.4cm}p{5.1cm}@{}}
\toprule
\textbf{Model / family} & \textbf{Type} & \textbf{Objective / data} & \textbf{Transfer value to animals} \\
\midrule
CPC / wav2vec / wav2vec~2.0~\cite{vandenoord2018cpc,schneider2019wav2vec,baevski2020wav2vec2} & SSL & contrastive / masked, raw audio & foundation bark encoder from unlabeled audio \\
HuBERT / WavLM / data2vec~\cite{hsu2021hubert,chen2022wavlm,baevski2022data2vec,baevski2023data2vec2} & SSL & masked unit prediction & robust, denoised vocalization features \\
wav2vec-U~\cite{baevski2021wav2vecu} & SSL & transcription-free & learn without a parallel corpus \\
XLS-R / MMS / USM~\cite{babu2022xlsr,pratap2023mms,zhang2023usm} & SSL & 100s--1000s of languages & one encoder across species \\
CTC / LAS / Conformer~\cite{graves2006ctc,chan2016las,gulati2020conformer} & ASR & alignment-free / attention & map variable-length signal to labels \\
Whisper~\cite{radford2023whisper} & ASR & 680k h weak supervision & robust off-the-shelf acoustic encoder \\
LibriSpeech / Common Voice / SUPERB~\cite{panayotov2015librispeech,ardila2020commonvoice,yang2021superb} & data & corpora + benchmarks & template for a ``bark corpus'' \\
SoundStream / EnCodec~\cite{zeghidour2021soundstream,defossez2022encodec} & codec & residual VQ & discretize vocalizations into tokens \\
AudioLM / VALL-E / UniAudio~\cite{borsos2023audiolm,wang2023valle,yang2024uniaudio} & codec-LM & token language modeling & search for sequence structure \\
AudioPaLM / SALMONN / Qwen-Audio~\cite{rubenstein2023audiopalm,tang2024salmonn,chu2024qwenaudio} & audio-LLM & audio + text fusion & bind acoustic signal to language \\
AVES / NatureLM / Perch / BEANS~\cite{hagiwara2023aves,robinson2024naturelm,ghani2023perch,hagiwara2023beans} & bioacoustic & animal audio SSL & cross-species transfer, already realized \\
\bottomrule
\end{tabular}
\end{table}